\section{System Design}
\label{sec:design}

\system selects which Agents' histories deserve protection (O1),
which KV regions within those histories should be retained (O2), and
how much HBM replication may occupy as foreground demand changes (O3).

Figure~\ref{fig:system-overview} shows four cooperating components.
\emph{Tracker} estimates remaining workflow work and recovery urgency.
\emph{Router} chooses execution workers and replica destinations.
\emph{KV Master} selects Agent histories and KV regions, evaluates
candidate placements, and manages replica retention and reclamation.
\emph{Adapter} adjusts the HBM replication budget to foreground demand. Rank-local
Connectors execute transfers and report completion to KV Master.

\begin{figure}[t]
 \centering
 \includegraphics[width=\columnwidth]{figures/system-overview.pdf}
 \caption{Four-component architecture. Each DP worker runs TP=2 and
 has execution HBM, replica HBM, and host DRAM within its own host
 boundary. The two workers are on different hosts in this example.
 Rank-matched Connector links carry copies between them; each worker's
 DRAM connects to its own Connectors. Blue paths carry KV and dashed
 paths carry control information. Router defines the peer destinations
 and memory choices in Section~\ref{sec:router}.}
 \label{fig:system-overview}
 \Description{Application, Controller, and Engine bands. Black callouts
 identify Tracker (1), Router (2), KV Master (3), and Adapter (4).
 Agent invoke
 A(1) forks into B(1) and C(1), which join at A(2). Tracker provides
 workflow estimates, Router exchanges candidate destinations with KV
 Master, and Adapter provides quotas. DP0 and DP1 each have TP0 and TP1
 HBM regions, rank-local Connectors, and their own DRAM allocation
 inside separate host boundaries. There is no standalone DRAM service.}
\end{figure}

\subsection{Tracker}
\label{sec:tracker}

Tracker maintains the residual invocation DAG and estimates remaining
work. Completed nodes retain history-consumer links but contribute no
execution; running nodes contribute only unfinished work.

\paragraph{One node is one invocation.}
A DAG node $v=(w,a,r)$ identifies the $r$th invocation of Agent $a$ in
workflow $w$. Thus, $a^{(1)}$ and $a^{(2)}$ are distinct nodes even when they
execute consecutively and use the same model. The later invocation may
extend the earlier input and accepted outputs. Reuse requires materialized state for the exact common token prefix
and compatible model, positions, format, and KV layout. Agent identity
or a dependency edge alone does not establish KV compatibility.

Successive invocations of an Agent may reuse KV from earlier
invocations (Figure~\ref{fig:invocation-reuse}). \system accounts
for these history dependencies when selecting and retaining KV replicas
for recovery. Tracker records history consumers, so selected KV produced
by $a^{(1)}$ can remain available for recovery in $a^{(2)}$ and $a^{(3)}$
without allocating a new backup for each call. Completing $a^{(1)}$
updates its remaining consumers and value; it does not automatically
retire its replicas. Tracker does not invent unseen future calls or
assign them guaranteed value. KV Master retains selected replicas of idle
histories under the retention policy in Section~\ref{sec:master-online}.

\begin{figure}[t]
 \centering
 \includegraphics[width=\columnwidth]{figures/invocation-reuse.pdf}
 \caption{Successive invocations of one Agent reuse growing history.
 In this append-only example, A(2) consumes the compatible prefix from
 A(1), and A(3) consumes the history extended by A(2). Each box in the top
 row is a separate DAG node. Lower bars show logical token histories,
 not complete backup allocations; only selected KV regions are retained.}
 \label{fig:invocation-reuse}
 \Description{Three successive invocation nodes A(1), A(2), and A(3) for
 Agent A. A(2)'s input contains history H1 plus new input U2; A(3)'s input
 contains H1, the extension E2 produced by A(2), and new input U3.
 H1 and H1 plus E2 are the compatible reusable prefixes.}
\end{figure}

\paragraph{Estimating urgency.}
A GBDT estimates remaining invocation duration from currently available
features: model and Agent role, input length, decoded-token progress,
elapsed execution, and completed calls' lengths and timings. Training
and calibration separate whole workflows. Forward and backward
critical-path passes~\cite{kelley1959cpm} produce earliest and latest
finishes, $EF_v$ and $LF_v$, and guarded slack
\begin{equation}
 s_v=\max\{0,LF_v-EF_v-g_v\},
 \label{eq:slack}
\end{equation}
where $g_v$ is a calibrated upper quantile of slack overestimation.
Successive invocations share path slack; it cannot be spent independently
at every node. Stale or invalid estimates grant no positive slack.
Slack orders urgent consumers, while KV Master's joint schedule evaluates
the effect of protecting their histories.

\subsection{Router}
\label{sec:router}

Router schedules invocations by queue load and compatible prefix
locality. Moving an invocation must justify its transfer and
reconstruction cost. Router also chooses where selected replicas reside.

\paragraph{Choosing a peer worker.}
For a history selected by KV Master, Router chooses a compatible
\emph{peer worker} $p(a)$ outside the execution worker's protected
failure domain. This peer is the preferred takeover target. Its GPU
memory used for replicas is \emph{Peer HBM}; the host DRAM allocated to
those replicas is \emph{Peer DRAM}. Both belong to the selected peer's
serving resources. Recovery on that peer consumes resident HBM directly
and transfers its DRAM-resident KV into the GPU when needed.

Peer choice considers projected takeover load, available memory, and
copy-path contention. An existing feasible binding is retained on ties.
Router fixes the proposed bindings for a planning snapshot so that KV
Master compares candidates under the same destinations and resource
conditions. Choosing a peer does not allocate a complete checkpoint;
only selected regions receive reservations and copies.

Peer DRAM is accounted at its actual host. If several workers share a
host, their DRAM reservations draw from the same physical capacity
ledger. HBM limits are checked separately for each TP rank. A different
worker on the same host is not a valid peer for host-failure protection.
The workflow history and replica directory also remain outside the
protected source failure domain.

\paragraph{Choosing HBM or DRAM.}
KV Master supplies a selected region set $S_a$ and its current marginal
priorities. A region identifies one layer and one historical token block
together with its required TP shards and attention metadata. Router
maps it to a memory tier on $p(a)$. The stable backup placement is
\begin{equation}
 S_a=H_a\mathbin{\dot\cup}D_a,\qquad
 H_a\subseteq\{b\in S_a:\operatorname{layer}(b)\le L\},
 \label{eq:placement}
\end{equation}
where $H_a$ resides in Peer HBM and $D_a$ in Peer DRAM. The disjoint
union denotes one managed backup per selected region; source execution
KV is accounted separately. Neither tier nor their union is required
to cover the full history.

The first $L$ model layers are eligible for HBM replication, with $L$
configured for a compatible model and layout. Router first offers the
selected regions in these layers to Peer HBM, ordered by KV Master's
profiled incremental workflow benefit of HBM over DRAM per byte.
Admissions must fit every required TP rank. Early-layer regions that
do not fit go to Peer DRAM, together with selected regions beyond $L$.
Thus, the depth limits placement, rather than requiring full copies
of the first $L$ layers. If the peer's DRAM budget is also insufficient,
Router returns the capacity shortfall to KV Master for reselection.

\paragraph{Maintaining the binding.}
On confirmed execution failure, Router fences the old epoch and
activates the peer. If a later invocation starts executing on its former
peer, Router chooses a new peer outside the new execution failure
domain. The existing local state remains useful for execution, but no
longer counts as an independent backup there. KV Master creates the
new selected replicas and publishes their protection only when copies
complete. Changing a binding accounts for migration, transient memory,
and any interval without a surviving backup. When the preferred peer
is unavailable, Router uses another compatible worker and charges
remote retrieval or reconstruction instead of assuming local hits.

\subsection{KV Master}
\label{sec:master}

KV Master translates Tracker's estimates into Agent selection (O1)
and region selection (O2), then evaluates Router's proposed placement.
Connectors execute the accepted plan and return copy status to Master.

\subsubsection{Selecting Agent Histories (O1)}
\label{sec:agent-protection}

\paragraph{Reuse and recovery value.}
For a selected set and placement $P$, let $F_w^0(P)$ be the predicted
finish of workflow $w$ during normal execution, and $F_w^f(P)$ its finish
after a common failure of source worker $f$. Normal execution includes
the Router's actual prefix locations, expected execution-KV retention,
and replica-copy interference. A backup earns no normal-reuse benefit
when the same prefix remains available in execution HBM. If the Router
uses a backup for a later invocation, any required transfer and
missing-state reconstruction precede that invocation.

The failure schedule includes all interrupted invocations and later
visible consumers of history lost with the worker. Restoration,
dependency reconstruction, and accepted-token replay precede unfinished
execution. Takeover admission and transfers contend on the resources
they actually share. A fixed ready-time rule, using guarded slack and
request ID to break ties, applies to every candidate. Master compares
\begin{equation}
 J(P)=(1-\lambda)\sum_w F_w^0(P)
       +\lambda\sum_f\pi_f\sum_w F_w^f(P),
 \label{eq:objective}
\end{equation}
where $0\le\lambda\le1$ is a configured recovery weight and $\pi_f$ are
fixed normalized source-failure weights, equal by default. These are
policy weights, not measured fault probabilities. The comparison fixes
the workload snapshot, Router targets, weights, and queue rules.

A single physical history can accelerate several later invocations,
but each invocation appears only once in each schedule. The objective
therefore avoids adding a full-history benefit for every reference.
Normal prefix reuse is evaluated separately from the fault-time slack
of the currently running call. At a join, the maximum predecessor
finish determines progress, preserving O1's non-additive effect.

\paragraph{Generating Agent candidates.}
Master first identifies Agent histories with active or visible future
consumers. Candidates include single-history changes and small bundles
covering co-bottleneck paths into an exposed join. Protecting only one
of two equally delayed branches may have zero effect, while protecting
both can advance the join. Each Agent candidate passes to O2 to
select useful regions; selection does not reserve a full checkpoint.
Master evaluates the resulting candidate, placement, and victims
together. Bounded search can miss other combinations and does not claim
global optimality. Without a valid DAG profile, it uses profiled
invocation time saved per byte with deterministic ties, retaining the
same correctness and capacity rules.

\subsubsection{Selecting KV Regions (O2)}
\label{sec:region-selection}

Within a candidate Agent history, Master selects sealed, materialized
regions whose retention reduces exposed waiting or reconstruction.
The logical set $S_a$ is formed before Router assigns its locations.
Master profiles the resulting placement and revises selection when its
cost or capacity is unsuitable; neither memory tier receives a full
history as a prerequisite.

\paragraph{Selecting KV regions.}
For region $b$ in layer $\ell$, recent completed native attention
observations supply
\begin{equation}
 \widehat p_{a,\ell,b}=
 \frac{\text{eligible selection units that selected }b}
      {\text{selection units for which }b\text{ was visible}}.
 \label{eq:demand}
\end{equation}
The unit follows the model's native selection rule; execution tiling
does not redefine it. Head duplicates count once per unit, and causal
visibility is determined before pruning. Small denominators use a
calibrated layer prior. Comparable visible blocks in dense attention
have access factor one. Frequency orders region candidates; it neither
guarantees future demand nor establishes a time saving.

Master considers individual regions and progressively larger hot sets,
including dependency bundles needed to make a proposed restoration
useful. Recent consecutive demand windows preserve correlated access
and native reveal order. Their size and the profiling work budget are
fixed parameters. Invalid sparse statistics use a conservative full
causal demand envelope for profiling, without forcing a full backup.
Neither selection nor placement changes native attention semantics.

\paragraph{Cost of a missing region.}
The profiler distinguishes a valid local execution hit, a Peer HBM hit, a
Peer DRAM fetch, and a region with no surviving copy. For the last case,
the engine reconstructs the required dependency closure from durable
tokens and reusable state. Rebuilding a block's upper-layer KV may
require lower-layer activations and causally preceding KV; the cost is
not an independent one-block Prefill. Native selector and positional
state are also dependencies. Shared reconstruction tasks are executed
and charged once within a recovery schedule.

If the engine cannot safely combine partial state with reconstruction,
it prefills the complete required prefix and gives those unusable
partial replicas no avoided-compute credit. A partial backup therefore
does not promise to avoid full reconstruction. Predictions can affect
latency, but every native demand must be satisfied before attention
continues. This is the correctness path for both later invocations
and failure recovery.

\paragraph{Measuring exposed delay.}
The profile schedules copies, bounded prefetches, dependency computation,
and unfinished invocation work on a common timeline. Replica creation
uses the inter-worker path; restoration from the selected peer's DRAM
uses its host-to-device path. Concurrent work shares the corresponding
network, host-memory, and GPU resources. A demand-driven
fetch cannot begin before the selector reveals it. If $q_\ell$ is the
time a layer is ready to consume history, $a_\ell(P)$ the time its last
required region becomes available, and $c_\ell$ its remaining compute,
then
\begin{equation}
 \begin{aligned}
 e_\ell(P)&=[a_\ell(P)-q_\ell]^+,\\
 f_\ell(P)&=q_\ell+e_\ell(P)+c_\ell.
 \end{aligned}
 \label{eq:overlap}
\end{equation}
Availability may require a transfer or reconstruction. Query and
selector preparation, resource queues, and dependencies determine
readiness; later layers depend on earlier finishes. Only computation
that remains to execute can hide a transfer. The resulting schedule
feeds Equation~\ref{eq:objective}; historical frequency is not
used as a substitute for exposed delay.

\begin{figure}[t]
 \centering
 \includegraphics[width=\columnwidth]{figures/selective-recovery.pdf}
 \caption{Partial selection before placement, on one selected peer, with $L=2$.
 Five of twelve logical regions are selected. Peer HBM holds L1/B0 and
 L1/B2; L2/B1 falls back to Peer DRAM because the illustrative HBM quota is
 full. Selected L3/B0 and L3/B2 also go to Peer DRAM. Unselected cells have
 no managed backup. A demand uses a surviving copy or reconstructs its
 missing dependencies. Each cell denotes all required TP shards;
 this is a constructed state, not a timing result.}
 \label{fig:selective-recovery}
 \Description{A selected-region mask followed by an exclusive Peer HBM/Peer DRAM
 placement map. Two selected first-layer regions reside in Peer HBM, an
 overflowing second-layer region resides in Peer DRAM, and two third-layer
 regions reside in Peer DRAM. Seven regions are unselected. The two tiers
 together store five regions, not a complete twelve-region checkpoint.}
\end{figure}

\subsubsection{Online Management}
\label{sec:master-online}
\label{sec:execution}

A candidate includes both its selected regions and any victims needed
to make Router's placement fit. Master considers removals in increasing
recomputed workflow loss per released byte, then evaluates the complete
replacement. Optional changes require $J(P)-J(P')>\epsilon$, a configured
hysteresis margin, and feasible transient memory and copy reservations.
Mandatory capacity relief can instead require deselection and its
explicit reconstruction fallback. Work and copy-byte limits bound
planning overhead and churn; no candidate claims globally optimal
selection or placement.

Planning runs on invocation arrival and completion, history or DAG
updates, memory pressure, and a fixed block-aligned interval $Z$ of
accepted tokens. Prefill and decode offer immutable materialized regions
as they become available. Final invocation completion offers its sealed
state and updates future-consumer references. The interval controls
review frequency; it is not a complete-checkpoint or replay-tail bound.

Algorithm~\ref{alg:protection} applies O1 and O2 before asking Router to assign
memory locations. Every candidate comparison includes current replicas,
victims, migrations, and normal-service copy interference. Version checks
cover the DAG, history, manifest, and quotas. A stale or infeasible
plan is requeued without advertising its proposed replicas. Candidate
gain is normalized by changed destination bytes for ordering; each
tier's physical capacity is enforced separately. Shared host DRAM is
charged once across all of that host's worker allocations.

\par\medskip\noindent\begin{minipage}{\columnwidth}
\raggedright
\refstepcounter{designalgorithm}\label{alg:protection}
\textbf{Algorithm \thedesignalgorithm: Select, place, and reconcile}
\par\smallskip
\begin{tabularx}{\columnwidth}{@{}r@{\quad}X@{}}
\toprule
1 & Reconcile copies; snapshot invocations and history references \\
2 & $P\gets$ current placement; apply mandatory capacity relief \\
3 & $A\gets\textsc{AgentCandidates}(P)$, including joint bundles \\
4 & \textbf{while} planning work budget remains \textbf{do} \\
5 & \quad Form candidate sets $S_c\gets\textsc{SelectRegions}(A,P)$ \\
6 & \quad For each $c$, get $P_c\gets\textsc{Router.Place}(S_c)$ \\
7 & \quad Fit victims and re-place if a budget is exceeded \\
8 & \quad Reprofile each whole plan, including missing-state work \\
9 & \quad $P^*\gets$ best feasible gain/byte candidate \\
10 & \quad \textbf{if} no gain exceeds $\epsilon$ \textbf{then break} \\
11 & \quad $P\gets P^*$; refresh affected candidates \\
12 & \textbf{if not} \textsc{ReserveIfCurrent}$(P)$ \textbf{then} \\
13 & \quad Requeue event; \textbf{return} \\
14 & Copy selected regions directly to their assigned tier \\
15 & Publish completed objects; retire old owners after drain \\
\bottomrule
\end{tabularx}
\end{minipage}\par\medskip

\paragraph{Publishing and retaining partial state.}
A manifest records each committed region's token interval, layer,
required shards, compatibility metadata, and current owner. Master
publishes a region only after Connectors validate all required objects;
pending copies cannot satisfy a lookup. For migration, old and new
allocations are both charged until the destination commits and old
readers drain. Failed copies leave the previous committed owner valid
when it survives. Retention follows live history references and current
selection value, including successors of a completed invocation.

The highest saved token position is not a complete checkpoint frontier.
Only an explicitly verified all-layer, all-shard prefix may be treated
as complete. Sparse manifests preserve holes, and the planner includes
their reconstruction cost. References from future invocations inform
value without permanently pinning replicas; only physical readers and transfers delay release.
With no active or visible future consumer, a configured idle-history
window permits retention subject to pressure. The history expires
at that limit; workflow completion releases its remaining logical references. Physical reuse
by another compatible consumer or an existing reader delays release.

\paragraph{Restoring a partial history.}
For a subsequent or interrupted invocation, Master leases the compatible
surviving regions and Connectors supply the required transfers. The
engine combines execution hits, Peer HBM hits, Peer DRAM reads, and
dependency reconstruction without omitting native attention demands.
Only validated, materialized KV can serve a prefix. Accepted tokens
are replayed where necessary without emitting them twice, extending
checkpoint-and-replay recovery~\cite{strati2024dejavu,cao2026lumen}.
Full-prefix Prefill remains the fallback when partial restoration is
unsupported or no useful state survives. Prediction changes latency;
compatibility, committed copies, and complete dependency evaluation
determine which state may safely be consumed.

\subsection{Adapter}
\label{sec:capacity}

Adapter realizes O3 by determining the HBM replication capacity
foreground execution can support. For takeover rank $k$,
$B^{\rm pool}_k$ excludes weights and fixed allocations,
$B^{\rm serve}_k$ covers the serving high watermark and execution
reservations, and $B^{\rm reserve}_k$ reserves growth and workspace.
The ceiling is
\begin{equation}
 C_k=\max\{0,\min(B^{\rm cap}_k,
 B^{\rm pool}_k-B^{\rm serve}_k-B^{\rm reserve}_k)\}.
 \label{eq:ceiling}
\end{equation}
HBM replica use is not part of $B^{\rm serve}_k$ and is not subtracted twice.
Committed, reserved, copying, and retiring allocations all count toward
physical occupancy. Distinct TP ranks have distinct ceilings.

Adapter uses additive increase and multiplicative
decrease~\cite{chiu1989aimd}. Allocation stalls, reserve violations, or
memory-related admission waits multiply quota $B_k$ by $0<\gamma<1$.
After $m$ stable windows below a lower wait threshold, Adapter adds
$\Delta$ bytes if Master reports beneficial early-layer candidates
currently in Peer DRAM or deferred for HBM capacity. The signal is computed
before quota filtering, allowing growth from zero. Updates are
page-aligned and clipped to $C_k$; missing telemetry disables growth.
Adapter changes total quota, while the configured depth $L$ remains
unchanged.

When quota shrinks, Master stops new HBM admissions and asks Router
to move the least valuable HBM placements to the same peer's DRAM if
their regions remain selected. The plan reserves that host's DRAM
space; Connectors copy and validate the destination. Master then
redirects the manifest and releases the HBM allocation after readers
drain. If that peer's DRAM cannot accommodate the migration, Master
reselects the finite backup set, explicitly removing regions and their
predicted protection benefit as needed.
Their token histories still permit reconstruction. Mandatory pressure
relief need not improve $J$, but cannot bypass active physical leases.
Foreground admission uses actual occupancy; a reduced quota alone
does not free memory. Pending copies may be cancelled, and allocations
wait when immediate physical space is unavailable.

The four components specify a design. Integrated measurements are
still needed to establish recovery gains, cross-invocation reuse, and
normal-service overhead.
